% Reciprocal proofs for the source-specific 1987 tutorial route.
\section{Klop 1987: abstract, residual, and strategy arguments}

\begin{annexitem}{annex:p12-klop1987-ne-normalization}{thm:p12-klop1987-ne-normalization}{Projection of an infinite reduction through a non-erasing step}
\begin{proof}
SN implies WN by well-founded induction on reduction: a term with no
successor is normal, and otherwise a successor has a normal-form descendant
by induction. For the converse we prove the stronger preservation fact:
in a regular NE system, if \(t\) admits an infinite reduction and
\(t\to s\), then \(s\) also admits an infinite reduction.

Fix the contracted occurrence \(u\) in \(t\) and give it an origin
label. Along an infinite reduction
\(t=t_0\to t_1\to t_2\to\cdots\), keep its residual set
\(U_i\). Initially \(U_0=\{u\}\). Each residual set is finite
and disjoint. Duplication of an enclosing variable instance copies its
label to disjoint leaves; contraction of a labelled residual removes it;
all other labels survive by the same residual rules as in
Lemma~\ref{thm:p12-klop1987-developments}. Non-overlap prevents an
unlabelled step from destroying a labelled operator pattern. Non-erasure
prevents destruction by erasing the variable instance containing it.
Let \(s_i\) be the endpoint obtained by developing \(U_i\) in
\(t_i\); then \(s_0=s\).

Commute each step \(t_i\to t_{i+1}\) past this development.
If it contracts a labelled residual, the canonical-development endpoint
does not change, so \(s_i=s_{i+1}\). If it is disjoint from the
labelled set, its unchanged occurrence contracts in \(s_i\), giving
\(s_i\to s_{i+1}\). If it is above labelled occurrences, all of
those lie in its variable instances. Developing them first gives a new
matching substitution; contracting the same outer rule then gives exactly
the right side with the selected residuals developed. Thus there is again
a genuine step \(s_i\to s_{i+1}\). If it is below a labelled
occurrence, it lies inside one of that occurrence's variable instances.
Developing the labelled root copies that inner step to every right-side
occurrence of the variable. NE guarantees at least one such occurrence,
and finitely many disjoint copies give a positive finite reduction
\(s_i\to^+s_{i+1}\). Other disjoint labelled roots do not affect
it. These are all relative position cases. Identity-valued rules still
give genuine counted steps; syntactic equality of endpoints does not
make a nonempty sequence empty.

There are infinitely many steps not contracting a labelled residual.
Otherwise, after the last such step, the remaining infinite tail would
be a development of the fixed finite residual set at that stage, contrary
to finite-development termination. Each of the infinitely many unlabelled
steps projects to a positive finite reduction of the \(s_i\), while
labelled steps project to an empty reduction. Concatenating the projected
finite segments therefore gives an infinite reduction from \(s\).
This proves the preservation fact.

If a term with an infinite reduction had a finite reduction to a normal
form, apply the preservation fact at every step of that finite reduction.
Its normal-form endpoint would then admit an infinite reduction, which
is impossible. Therefore WN implies SN locally. Quantifying the local
equivalence over all terms gives the system-wide version. The argument
uses NE precisely to keep the projected inner segment nonempty; without
it the erasing \(K\)-example makes the assertion false.
\end{proof}
\end{annexitem}

\begin{annexitem}{annex:p12-klop1987-ars-core}{thm:p12-klop1987-ars-core}{Core ARS implications with the countability boundary}
\begin{proof}
First suppose CR. Any finite conversion can be joined: induct on its
length, keeping a common descendant of the first endpoint and the current
endpoint. A forward next step and the current joining reduction are a
confluent peak at the current endpoint; join them. A backward next step
can simply be followed by the existing joining reduction. This produces
a join of the conversion endpoints. Conversely the endpoints of a peak
are convertible through its source, so conversion joinability implies CR.
The mixed peak condition implies the same conversion-joining induction;
CR plainly supplies that condition. The conditions on the reflexive
transitive relation in Proposition 1.2 coincide because that relation is
already reflexive and transitive.

If a normal form is convertible with \(b\), conversion joinability
provides a common descendant. The normal form has no step, so that
descendant is the normal form itself; this proves NF. If two normal forms
descend from \(a\), they are convertible; NF makes one reduce to the
other, and irreducibility forces identity. Hence NF implies UN.

Assume UN and WN. For a peak from \(a\) with endpoints \(b,c\),
choose normal forms \(n_b,n_c\) below those endpoints. They are both
below \(a\), so UN identifies them. This proves CR. For any reduction
sequence beginning at \(a\), choose a normal form of \(a\).
Every sequence member has some normal form by WN; it is also a normal
form of \(a\), and UN identifies it with the chosen one. Thus this one
normal form is below every sequence member, proving Ind.

If Ind and Inc held for an infinite sequence \(a_0\to a_1\to\cdots\),
Ind would provide \(b\) below every \(a_i\). Inc would imply
\(|a_i|\leq|b|\) for every \(i\), while the strictly increasing
natural-number sequence \(|a_i|\) is unbounded. This contradiction
proves SN.

For CP implies CR, take a peak from \(a\) ending at \(b,c\).
Cofinality gives sequence members \(a_i,a_j\) below \(b,c\),
respectively. The later member of this same sequence is below both,
giving their join. No countability assumption is used in this direction.
Conversely suppose CR and that \(A\) is countable. Enumerate the
descendants of a fixed \(a\) as \(b_0,b_1,\ldots\), repeating
elements if the set is finite. Begin at \(a\). At stage \(i\),
join the current endpoint with \(b_i\), both descendants of \(a\),
and append the finite reduction from the endpoint to that join. Every
\(b_i\) reaches an endpoint occurring in the concatenated reduction.
If infinitely many actual steps are appended, this is an infinite cofinal
reduction. If only finitely many are appended, its final endpoint remains
a common descendant of every subsequently enumerated element, and the
finite reduction is cofinal. This construction establishes CP for each
\(a\). The enumeration is exactly where countability enters; it is
not supplied for an uncountable ARS by the argument.
\end{proof}
\end{annexitem}

\begin{annexitem}{annex:p12-klop1987-one-rule-puzzle}{thm:p12-klop1987-one-rule-puzzle}{Bottom-up normalization and the self-loop}
\begin{proof}
The term \(q=F(a,F(a,F(b,b)))\) matches the rule at its root with
\(x=a,y=F(b,b)\); its contractum is literally \(q\). Repeating
that genuine rewrite gives an infinite reduction, so the system is not SN.

For WN, normalize finite terms by induction on their tree structure.
Variables and constants have no redex. For \(F(s,t)\), use the
induction hypotheses to reduce its children to normal forms \(s',t'\).
Context closure gives \(F(s,t)\to^*F(s',t')\). If this root is
not a redex, both children are normal, so the entire term is normal.
Otherwise \(s'=a\) and \(t'=F(v,w)\), with normal \(v,w\).
The root contracts to \(F(v,F(v,F(b,b)))\).

If \(v\ne a\) syntactically, none of the new \(F\)-nodes has
left child \(a\), except nodes already inside the normal term \(v\).
There are no redexes there either. Thus this result is normal.
If \(v=a\), the result is the loop term \(q\), but there is a
different finite reduction: contract its proper subterm
\(F(a,F(b,b))\) to \(F(b,F(b,F(b,b)))\), then contract the
remaining outer root. Its final result is
\(F(b,F(b,F(b,b)))\), which is normal because all its \(F\)-nodes
have left child \(b\), not \(a\).
This proves existence of a normal form in every induction case.
It does not prohibit the root self-loop or assert that every reduction
uses this bottom-up choice.
\end{proof}
\end{annexitem}

\begin{annexitem}{annex:p12-klop1987-semantic-un}{thm:p12-klop1987-semantic-un}{From equational soundness to syntactic separation}
\begin{proof}
Evaluation of a rule instance under any assignment gives equal values,
since evaluating a substituted variable defines another assignment to
the variables of the original rule. Equal values remain equal after
evaluation of an outer context, by induction on that context. Thus every
rewrite step is sound in the algebra, and equality of evaluations is
preserved along every finite reduction. If \(t\) reduces to normal
forms \(n,m\), the algebra satisfies \(t=n\) and \(t=m\)
under every assignment, hence satisfies \(n=m\). The stipulated
separation of syntactically distinct normal forms gives \(n=m\)
as terms. This is UN and uses neither WN nor SN nor groundness.
\end{proof}
\end{annexitem}

\begin{annexitem}{annex:p12-klop1987-regular-confluence}{thm:p12-klop1987-regular-confluence}{All positions in the parallel diamond}
\begin{proof}
Represent each parallel contraction by its finite set of disjoint
positions. Prove the diamond by induction on the finite source term.
If neither set includes the root, both sets split into disjoint sets
inside the immediate children. Apply the induction assertion to the
two child contractions in each child, including empty contractions.
Assemble the common child reducts under the unchanged root operator.
The union of the child joining sets is disjoint, and hence gives a
parallel join on each side. For a variable or constant with no selected
root this is the empty contraction case.

If both sets include the root, disjointness says each consists of that
root alone. Non-ambiguity excludes two distinct rule patterns matching
there. A left-linear pattern also determines its matching assignment
uniquely from the concrete input at its variable leaves. Consequently
the two contractions have the same result, and join by empty sets.

Suppose exactly the first set includes the root. It is a singleton root
contraction \(l^\sigma\to r^\sigma\). Every position in the other
set must lie in an instance of a variable of \(l\): otherwise an inner
left-hand-side pattern would overlap a nonvariable position of \(l\),
contrary to regularity. Since \(l\) is left-linear, each variable has
one source occurrence, so the selected inner contractions define a
single substitution \(\sigma'\) with
\(\sigma(x)\to_{\parallel}\sigma'(x)\) for each variable
\(x\) of \(l\). Unaffected variables have empty contractions.
The inner-first endpoint is exactly \(l^{\sigma'}\), not an
inconsistent assignment to repeated variable leaves. Contract its root
to \(r^{\sigma'}\).

The root-first endpoint \(r^\sigma\) reaches \(r^{\sigma'}\)
by copying the selected inner set into every occurrence of its variable
on the right. Distinct variable leaves of the finite right side are
disjoint, and within each copy the selected positions are disjoint.
Thus their union is a disjoint parallel contraction. Duplication gives
more copies; erasure gives no copy and an empty arm. Both are legitimate.
This proves the remaining case; exchanging the two sets proves its
symmetric form. These cases exhaust root and nonroot positions.

The diamond property gives confluence of the reflexive transitive
parallel relation by the finite grid argument: commute the first step
of one path successively past every step of the other, and repeat for
the remaining steps of the first finite path. Each cell is a parallel
diamond; the far corner joins the two endpoints. Each ordinary step is
a singleton parallel step, and each parallel step sequentializes into
ordinary steps. The two reflexive transitive relations therefore agree.
Their confluence proves ordinary CR without termination.
\end{proof}
\end{annexitem}

\begin{annexitem}{annex:p12-klop1987-cofinal-normalizing}{thm:p12-klop1987-cofinal-normalizing}{Cofinality and the role of irreducibility}
\begin{proof}
Let \(n\) be a normal form below \(t\). Cofinality of the strategy
orbit gives an iterate \(\mathbb F^k(t)\) below \(n\).
A normal form admits only the empty finite reduction, so that iterate
is literally \(n\). Thus the strategy normalizes every normalizable
input. If CR holds, two normal forms of \(t\) have a common descendant;
irreducibility identifies each with that descendant, giving uniqueness.
The definitions impose no obligation on merely normalizing strategies
to overtake a reduct of an input that has no normal form. The source's
\(\Omega\) example exploits exactly this additional cofinality demand.
\end{proof}
\end{annexitem}

\begin{annexitem}{annex:p12-klop1987-developments}{thm:p12-klop1987-developments}{Residual labels, termination, and the canonical development}
\begin{proof}
Fix a finite input term and a finite selected set \(U\) of its redex
occurrences. Give each selected occurrence a distinct origin label.
Across a contraction of a selected redex, unselected positions disjoint
from it survive unchanged. A selected redex below it lies in a variable
instance by non-overlap. Its label is copied at every occurrence of that
variable on the right, or erased if the variable is absent there.
A selected ancestor retains its label because rewriting inside one of
its variable instances preserves its left-linear pattern match. No new
label is attached to a newly created, unselected redex. These rules define
the residual selected set after each development step.

Order the finitely many original labels by their depths in the original
selected containment forest: depth zero for selected occurrences without
a selected ancestor, depth one for their nearest selected descendants,
and so on. An ancestor's origin depth is strictly smaller than a
descendant's. In a current term, an ancestor label can contain only
labels originating strictly below it. Contracting a label of depth
\(d\) decreases the number of depth-\(d\) residuals by one, leaves
all counts at smaller depths unchanged, and can change only counts at
larger depths through duplication or erasure. The finite tuple of
nonnegative counts, ordered lexicographically with shallower depths
first, strictly decreases. Lexicographic order on a fixed finite product
of \(\mathbb N\) is well-founded, by induction on the tuple length.
Hence every development terminates; it cannot be made infinite by
duplicating deeper labels. Each step leaves a finite term and finitely
many residual labels.

Define a canonical outcome recursively on the original marked term.
If the root is not selected, retain its operator and apply the construction
to its children. If the root is selected and has match \(l^\sigma\),
all deeper selected occurrences lie in the separate variable instances.
Recursively develop each such instance to \(\widehat\sigma(x)\),
then output \(r^{\widehat\sigma}\). This recursion is well-founded
on proper source subterms. Multiple occurrences on the right use the
same developed variable instance; erasure omits it.

The same canonical operation can be applied to any current labelled term.
It is invariant under a development step. For a step at a selected root,
developing that root first produces exactly the right side with the
residual labels copied or erased as specified; developing its variable
instances afterwards gives \(r^{\widehat\sigma}\), which is also
the recursive definition before that step. For a step below a selected
root, the step occurs within a variable instance. Induction on that
proper subterm leaves its canonical developed value unchanged, and
therefore leaves the enclosing right-side substitution unchanged.
For a step below an unselected root, apply the same induction to its
affected child and retain the other children. These cases establish
invariance for every labelled contraction, including arbitrary outer
contexts.

A complete development has no selected residuals, so its canonical
outcome is its current unlabelled term. Invariance therefore identifies
all complete-development endpoints with the canonical original outcome.
Termination guarantees that repeatedly choosing a residual reaches such
an endpoint. A zero-redex term is fixed. A non-normal finite term has
at least one selected redex when all its redexes are selected, and the
construction contracts at least one actual rule instance, giving a
positive ordinary reduction to the endpoint even if the endpoint happens
to be syntactically the input itself. For an effective rule presentation,
finite matching and the specified recursive construction compute that
endpoint. This establishes well-defined full substitution, but not yet
its cofinality on arbitrary newly created redexes.
\end{proof}
\end{annexitem}
